\chapter{The Bullshit Meter}
\label{chap:bullshit-meter}

% EPISODE COVERAGE: Episodes 6 and 7 revisited through elaboration cost,
% plus Episode 10's HeartbeatProcess as forward reference.
% KEY CONCEPT: The heartbeat count is the compiler measuring the logical
% distance between context and conclusion.

\section{What the Compiler Finds Expensive}
\label{se:what-compiler-finds-expensive}

% In Lean 4, every elaboration step consumes a heartbeat.
% The device raises maxHeartbeats repeatedly:
%
%   set_option maxHeartbeats 400000    -- Episode 7
%   set_option maxHeartbeats 2000000   -- Episode 13
%   set_option maxHeartbeats 4000000   -- Episode 15
%   set_option maxHeartbeats 0         -- Episode 15 final lines
%
% These are not programmer impatience. They are measurements.
% Each increase records the compiler's honest report of how hard it is
% working to verify the tower so far.
%
% Sidebar: What is elaboration? Filling in implicit arguments, resolving
% typeclasses, confirming types. The heartbeat count is the clock ticking
% while it works.

\begin{gphenom}{Pythagoras-Planck}
\label{ph:pythagoras-planck}
\GPhObs Planck's constant $h = 6.626 \times 10^{-34}$ J$\cdot$s sets a minimum
quantum of action. No physical process can transfer less than one quantum.
Below $h$, no distinction is admissible.
\GPhCon The ledger is finite at every scale. There exists a smallest admissible
distinction, below which no further elaboration is possible.
\GPhSeq The compiler's heartbeat is the formal analog of $h$: the minimum
unit of elaboration cost. Below one heartbeat, no verification step
is admissible. The tower cannot be built for free.
\end{gphenom}

\section{The Bullshit Meter Annotations}
\label{se:bullshit-meter-annotations}

% The comments in the code: -- Bullshit meter approximately N.
% Show a sampling from different episodes to illustrate the meter climbing.
%
% Early episodes: small numbers.
% Episode 6: climbs as arithmetic and splines are introduced.
% Episode 7: the instance cascade. The meter spikes.
% Episode 15: 32 typeclass parameters. The meter is pegged.
%
% Frankfurt's bullshit: "On Bullshit" (2005) distinguishes lies from bullshit.
% The bullshit meter measures the accumulated overhead of prior commitments
% that the compiler must now carry -- regardless of whether those commitments
% were well-chosen.
%
% Sidebar: Frankfurt's bullshit. The bullshitter does not care about truth.
% The bullshit meter cares about nothing but truth. It measures exactly how
% much truth the compiler has been forced to carry.

\section{Precision Has a Price}
\label{se:precision-has-a-price}

% The spline epsilon: the informational quantum.
% After threading through all observed points and satisfying all knot
% constraints, there remains a family of valid splines that cannot be
% distinguished by any admissible measurement within the current ledger.
%
% The compiler analogy: the gap between definitional and propositional equality.
% Below epsilon, no new distinction is admissible in either system.
%
% Connect forward: Chapter 8 (Stress) in the instrument develops the spline
% epsilon in full. Here we show its Lean precursor in FINITE_ELEPHANT.

\begin{gphenom}{von Neumann}
\label{ph:von-neumann}
\GPhObs Every floating-point system has a machine epsilon: the smallest
number that, added to 1.0, produces a result distinguishable from 1.0.
Below machine epsilon, the arithmetic system cannot record a difference.
\GPhCon The instrument cannot resolve distinctions smaller than its machine
epsilon. The ledger cannot record events below this threshold.
No refinement below this scale produces a new entry.
\GPhSeq \lean{FINITE\_ELEPHANT}'s \lean{finite?} predicate is the machine epsilon
test for the tower: below epsilon, the sequence has converged and no
further refinement is admissible. The compiler stops not because it
is tired but because there is nothing left to distinguish.
\end{gphenom}

\section{When Refinement Stops}
\label{se:when-refinement-stops}

% The Richardson Effect: measure a coastline at 100km resolution and you
% get one number. At 10km, a larger number. At 1km, larger still.
% The measured length diverges as resolution increases.
%
% This is not a paradox -- it is the correct behavior of any instrument
% near its resolution limit. At some scale, the instrument can no longer
% distinguish the signal from the noise.
%
% In the device, this stopping point is formalized by FINITE_ELEPHANT's
% finite? predicate.
%
% The name FINITE_ELEPHANT: the elephant in the room of finite element
% analysis is that all smooth fields are approximated by piecewise
% polynomial functions. FINITE_ELEPHANT is the honest acknowledgment
% that the elephant exists and is finite.

\begin{gphenom}{Richardson}
\label{ph:richardson}
\GPhObs Lewis Fry Richardson measured that the reported length of a national
border depends entirely on the length of the ruler used to measure it.
Shorter rulers find longer borders. The coastline has no definite length.
\GPhCon A fractal boundary has no well-defined length at infinite resolution.
The instrument must choose a resolution and accept the result as approximate.
Below the chosen scale, no further distinction between segments is admissible.
\GPhSeq The spline's informational quantum $\varepsilon$ is the compiler's
chosen resolution: the point below which no further distinction is admissible.
The ledger does not grow forever. It stops when refinement stops buying anything.
\end{gphenom}

% End of chapter: the compiler has been measured. Its elaboration cost is
% a real number, rising with each new layer. Its resolution limit is real.
% The instrument is calibrated. It knows its own limits.

